一面积为 \(A>0\) 的理想刚性真空环形干涉仪绕其面法向以恒定 \(\Omega\) 转动，周长为 \(L\)，线速度始终满足 \(|\Omega|\sqrt A\ll c\)。从同一旋转耦合点向两个相反方向注入频率 \(\omega_0=2\pi c/\lambda\) 的光，并在返回时比较相位，忽略色散和色散镜面延迟。
(1) 对半径 \(R\) 的圆环从实验室惯性系追及运动接收点的光程求双向返回时间及其一阶差。
(2) 将圆环结果推广为任意平面闭合路径，使用 \(\oint(\boldsymbol\Omega\times\boldsymbol r)\cdot d\boldsymbol\ell\) 推出只依赖有向面积的Sagnac差。
(3) 求两束的相位差及绕行 \(m\) 圈的增强因子，区分几何面积、面积矢量和波长定义。
(4) 若检测 \(N\gg1\) 个统计独立光子且干涉工作点最优，单次相位标准差的散粒下界为 \(1/\sqrt N\)，推出角速度灵敏度并讨论机械抖动、损耗及超相对论旋转的失效。
